\Needspace{12\baselineskip}
\subsection{Hash doble: comparar substrings rápido}\label{rec:hash}

\campo{Sirve para}{comparar dos pedazos de un string en O(1), calcular el
largo del prefijo común (LCP) y buscar patrones con búsqueda binaria.}
\campo{Entrada}{un string \codigo{s}, una cantidad de consultas \codigo{q} y
cada consulta con dos rangos \codigo{l1 r1 l2 r2}.}
\campo{Salida}{\codigo{SI} si \codigo{s[l1..r1)} es igual a
\codigo{s[l2..r2)}, o \codigo{NO} si no.}
\complejidad{Preprocesar O(n); cada comparación O(1); cada LCP O(log n).}

\lstinputlisting{codigo/hash.cpp}

\subsubsection*{Los rangos son medios abiertos}

\codigo{[l, r)} significa: incluye \codigo{l} y \textbf{no} incluye
\codigo{r}. Por eso \codigo{s[l..r)} se escribe \codigo{s.substr(l, r - l)}.
Un rango vacío es \codigo{l == r}.

\ojo{El hash es probabilístico: con dos módulos la probabilidad de error es
prácticamente nula, pero no es cero. Si el problema exige certeza absoluta,
usa un autómata de sufijos o compara los strings de verdad. No uses un solo
módulo si puedes evitarlo.}
